\newcommand{\inversion}[1]{\!\overleftrightarrow{\,#1\,}\!}

\newcommand{\reverse}{{\,\overleftrightarrow{\!(\cdots)\!}\,}}

\usepackage{tikz}

\newcommand{\crosssquare}{{\begin{tikzpicture}[scale=.15]
			\draw (0,0) -- (1,1) -- (1,0) -- (0,1) -- (0,0) -- (1,0) -- (1,1) -- (0,1) -- cycle;
\end{tikzpicture}}}

\newcommand{\cross}{{\begin{tikzpicture}[scale=.15]
			\draw (0,0) -- (1,1);
			\draw (1,0) -- (0,1);
\end{tikzpicture}}}

\newcommand{\triang}{{\begin{tikzpicture}[scale=.15]
			\draw (1,0) -- (0,0) -- (0,1) -- cycle;
\end{tikzpicture}}}

\newcommand{\trileg}{{\begin{tikzpicture}[scale=.15]
			\draw (1,1) -- (0,0);
			\draw (1,0) -- (0,0) -- (0,1);
\end{tikzpicture}}}

\newcommand{\cyclic}{{\circlearrowleft}_{ijk}}

\newcommand{\pdv}[2]{\frac{\partial {#1}}{\partial {#2}}}

\newcommand{\C}{\mathbb{C}}
\newcommand{\K}{\mathbb{K}}
\newcommand{\N}{\mathbb{N}}
\newcommand{\R}{\mathbb{R}}
\newcommand{\Z}{\mathbb{Z}}

\newcommand{\dilog}{\operatorname{Li_2}}

\newcommand{\cL}{\mathcal{L}}

\renewcommand{\d}{\mathrm{d}}

\renewcommand{\vec}{\boldsymbol}

{ \theoremstyle{definition}
	\newtheorem{Definition}{Definition}[section]
	\newtheorem{Example}[Definition]{Example}
	\newtheorem{ExampleH1}{Example: H1, part}
	\newtheorem{ExampleH2}{Example: H2, part}
	\newtheorem{ExampleH3}{Example: H3$_{\delta=0}$, part}
	\newtheorem{Remark}[Definition]{Remark}
}
\newtheorem{Lemma}[Definition]{Lemma}
\newtheorem{Theorem}[Definition]{Theorem}
\newtheorem{Proposition}[Definition]{Proposition}
\newtheorem{Corollary}[Definition]{Corollary}
\newtheorem{Assumption}[Definition]{Assumption}

\numberwithin{equation}{section}


\allowdisplaybreaks
\usepackage[a4paper,landscape,margin=20mm]{geometry}
\AddToHook{cmd/maketitle/after}{\markboth{Richardson and Vermeeren}{ABS multiform double-zero expansions}}
\pdfstringdefDisableCommands{\def\boldsymbol#1{#1}\def\bm#1{#1}\def\mathbb#1{#1}}
